% Options for packages loaded elsewhere
\PassOptionsToPackage{unicode}{hyperref}
\PassOptionsToPackage{hyphens}{url}
\PassOptionsToPackage{dvipsnames,svgnames,x11names}{xcolor}
\documentclass[
  spanish,
  11pt,
  a4paper,
  DIV=11,
  numbers=noendperiod]{scrartcl}
\usepackage{xcolor}
\usepackage[top=24mm,bottom=24mm,left=24mm,right=24mm,headheight=15pt,headsep=8mm]{geometry}
\usepackage{amsmath,amssymb}
\setcounter{secnumdepth}{5}
\usepackage{iftex}
\ifPDFTeX
  \usepackage[T1]{fontenc}
  \usepackage[utf8]{inputenc}
  \usepackage{textcomp} % provide euro and other symbols
\else % if luatex or xetex
  \usepackage{unicode-math} % this also loads fontspec
  \defaultfontfeatures{Scale=MatchLowercase}
  \defaultfontfeatures[\rmfamily]{Ligatures=TeX,Scale=1}
\fi
\usepackage{lmodern}
\ifPDFTeX\else
  % xetex/luatex font selection
\fi
% Use upquote if available, for straight quotes in verbatim environments
\IfFileExists{upquote.sty}{\usepackage{upquote}}{}
\IfFileExists{microtype.sty}{% use microtype if available
  \usepackage[]{microtype}
  \UseMicrotypeSet[protrusion]{basicmath} % disable protrusion for tt fonts
}{}
\makeatletter
\@ifundefined{KOMAClassName}{% if non-KOMA class
  \IfFileExists{parskip.sty}{%
    \usepackage{parskip}
  }{% else
    \setlength{\parindent}{0pt}
    \setlength{\parskip}{6pt plus 2pt minus 1pt}}
}{% if KOMA class
  \KOMAoptions{parskip=half}}
\makeatother
\usepackage{graphicx}
\makeatletter
% Set default figure placement to htbp
\def\fps@figure{htbp}
\makeatother
\ifLuaTeX
\usepackage[bidi=basic]{babel}
\else
\usepackage[bidi=default]{babel}
\fi
% get rid of language-specific shorthands (see #6817):
\let\LanguageShortHands\languageshorthands
\def\languageshorthands#1{}
\setlength{\emergencystretch}{3em} % prevent overfull lines
\providecommand{\tightlist}{%
  \setlength{\itemsep}{0pt}\setlength{\parskip}{0pt}}


\newcommand{\practicaPdfCoverImage}{../resources/practica2/images/teaser.jpg}
% PDF style for Practica 9.
% A modest university handout look: classical type, clear hierarchy,
% restrained color, and durable code/callout styling.

\usepackage{graphicx}
\usepackage[export]{adjustbox}
\usepackage{fontspec}
\usepackage{microtype}
\usepackage{scrlayer-scrpage}
\usepackage{tikz}
\usepackage{fvextra}
\usepackage{caption}
\usepackage{subcaption}
\usepackage[skins,breakable]{tcolorbox}

\providecommand{\practicaPdfCoverImage}{}

\definecolor{uzaInk}{HTML}{172033}
\definecolor{uzaMuted}{HTML}{5F6B7A}
\definecolor{uzaBlue}{HTML}{1F5F8B}
\definecolor{uzaBlueDark}{HTML}{153E5C}
\definecolor{uzaCyan}{HTML}{EAF4F7}
\definecolor{uzaLine}{HTML}{D7E2EA}
\definecolor{uzaCodeBg}{HTML}{F5F7F9}
\definecolor{uzaTip}{HTML}{237A57}
\definecolor{uzaWarn}{HTML}{A65F00}

\IfFontExistsTF{TeX Gyre Pagella}{
  \setmainfont{TeX Gyre Pagella}
}{}
\IfFontExistsTF{TeX Gyre Heros}{
  \setsansfont{TeX Gyre Heros}
}{}
\IfFontExistsTF{TeX Gyre Cursor}{
  \setmonofont{TeX Gyre Cursor}[Scale=0.86]
}{
  \setmonofont{Latin Modern Mono}[Scale=0.88]
}

\KOMAoptions{parskip=half}
\setlength{\parindent}{0pt}
\setlength{\parskip}{0.55em plus 0.15em minus 0.08em}
\linespread{1.035}

\addtokomafont{disposition}{\sffamily\color{uzaBlueDark}}
\addtokomafont{section}{\LARGE}
\addtokomafont{subsection}{\Large}
\addtokomafont{subsubsection}{\normalsize\bfseries}
\RedeclareSectionCommand[beforeskip=2.2em plus 0.4em minus 0.2em,afterskip=1.0em]{section}
\RedeclareSectionCommand[beforeskip=1.2em,afterskip=0.55em]{subsection}
\RedeclareSectionCommand[beforeskip=1.0em,afterskip=0.35em]{subsubsection}

\clearpairofpagestyles
\automark[section]{section}
\ihead{\sffamily\footnotesize\textcolor{uzaMuted}{Informática (30717)}}
\ohead{\sffamily\footnotesize\textcolor{uzaMuted}{\headmark}}
\cfoot{\sffamily\small\textcolor{uzaMuted}{\pagemark}}
\setkomafont{pageheadfoot}{\normalfont}
\KOMAoptions{headsepline=0.35pt}
\addtokomafont{headsepline}{\color{uzaLine}}

\makeatletter
\AtBeginDocument{%
  \hypersetup{
    linkcolor=uzaBlueDark,
    urlcolor=uzaBlue,
    citecolor=uzaBlueDark
  }%
  \@ifundefined{Highlighting}{%
    \DefineVerbatimEnvironment{Highlighting}{Verbatim}{
      commandchars=\\\{\},
      fontsize=\small,
      breaklines=true,
      breakanywhere=true
    }%
  }{%
    \RecustomVerbatimEnvironment{Highlighting}{Verbatim}{
      commandchars=\\\{\},
      fontsize=\small,
      breaklines=true,
      breakanywhere=true
    }%
  }%
  \@ifundefined{Shaded}{%
    \newenvironment{Shaded}{%
      \begin{tcolorbox}[
        enhanced,
        breakable,
        sharp corners,
        boxrule=0.25pt,
        colback=uzaCodeBg,
        colframe=uzaLine,
        borderline west={1.4pt}{0pt}{uzaBlue},
        left=2mm,
        right=2mm,
        top=1.4mm,
        bottom=1.4mm,
        before skip=0.75em,
        after skip=0.75em
      ]%
    }{%
      \end{tcolorbox}%
    }%
  }{%
    \renewenvironment{Shaded}{%
      \begin{tcolorbox}[
        enhanced,
        breakable,
        sharp corners,
        boxrule=0.25pt,
        colback=uzaCodeBg,
        colframe=uzaLine,
        borderline west={1.4pt}{0pt}{uzaBlue},
        left=2mm,
        right=2mm,
        top=1.4mm,
        bottom=1.4mm,
        before skip=0.75em,
        after skip=0.75em
      ]%
    }{%
      \end{tcolorbox}%
    }%
  }%
}
\makeatother

\newcommand{\PracticeCode}[1]{%
  \begin{Shaded}%
  \VerbatimInput[
    fontsize=\small,
    breaklines=true,
    breakanywhere=true
  ]{#1}%
  \end{Shaded}%
}

\makeatletter
\renewcommand{\maketitle}{%
  \begin{titlepage}
    \thispagestyle{empty}
    \begin{tikzpicture}[remember picture,overlay]
      \fill[uzaCyan] (current page.north west) rectangle ([yshift=-72mm]current page.north east);
      \fill[uzaBlue] (current page.north west) rectangle ([xshift=7mm]current page.south west);
    \end{tikzpicture}
    \vspace*{12mm}
    {\sffamily\small\bfseries\textcolor{uzaBlue}{INFORMÁTICA (30717)}}\par
    \vspace{5mm}
    {\sffamily\bfseries\fontsize{30}{35}\selectfont\textcolor{uzaInk}{\@title}}\par
    \vspace{3mm}
    {\sffamily\Large\textcolor{uzaMuted}{Grado en Estudios en Arquitectura}}\par
    \vspace{8mm}
    {\color{uzaBlue}\rule{\textwidth}{0.6pt}}\par
    \if\relax\detokenize\expandafter{\practicaPdfCoverImage}\relax
    \else
      \vspace{5mm}
      {\centering\includegraphics[width=\textwidth,height=78mm,keepaspectratio]{\practicaPdfCoverImage}\par}
    \fi
    \vfill
    {\sffamily\large\textcolor{uzaBlueDark}{Universidad de Zaragoza}}\par
    \vspace{2mm}
    {\sffamily\small\textcolor{uzaMuted}{Material de prácticas del curso 2025--2026}}\par
  \end{titlepage}
}
\makeatother

\captionsetup{
  font={small,sf},
  labelfont={bf,color=uzaBlueDark},
  textfont={color=uzaInk},
  justification=centering,
  singlelinecheck=false,
  skip=6pt
}
\captionsetup[subfigure]{
  font={small,sf},
  labelfont={bf,color=uzaBlueDark},
  justification=centering,
  singlelinecheck=true,
  skip=4pt
}

\newcommand{\PracticeCredit}[1]{\textcolor{uzaMuted}{#1}}
\newcommand{\PracticeImagePair}[6]{%
  \begin{figure}[tbp]
    \centering
    \begin{minipage}[t]{0.485\linewidth}
      \includegraphics[width=\linewidth]{#3}
      \vspace{2mm}
      \centering{\sffamily\footnotesize #4}
    \end{minipage}\hfill
    \begin{minipage}[t]{0.485\linewidth}
      \includegraphics[width=\linewidth]{#5}
      \vspace{2mm}
      \centering{\sffamily\footnotesize #6}
    \end{minipage}
    \caption{\label{#1}#2}
  \end{figure}
}

\renewcommand{\arraystretch}{1.16}
\setlength{\tabcolsep}{6pt}
\setkeys{Gin}{center}

% Quarto callout colors are referenced by generated tcolorbox options.
% Apply them at document start so they win over Quarto's defaults.
\AtBeginDocument{%
  \definecolor{quarto-callout-note-color}{HTML}{1F5F8B}%
  \definecolor{quarto-callout-note-color-frame}{HTML}{1F5F8B}%
  \definecolor{quarto-callout-tip-color}{HTML}{237A57}%
  \definecolor{quarto-callout-tip-color-frame}{HTML}{237A57}%
  \definecolor{quarto-callout-warning-color}{HTML}{A65F00}%
  \definecolor{quarto-callout-warning-color-frame}{HTML}{A65F00}%
  \definecolor{quarto-callout-important-color}{HTML}{9B2C2C}%
  \definecolor{quarto-callout-important-color-frame}{HTML}{9B2C2C}%
  \definecolor{quarto-callout-caution-color}{HTML}{B45309}%
  \definecolor{quarto-callout-caution-color-frame}{HTML}{B45309}%
}

\KOMAoption{captions}{tableheading}
\makeatletter
\AtBeginDocument{%
\ifdefined\contentsname
  \renewcommand*\contentsname{Tabla de contenidos}
\else
  \newcommand\contentsname{Tabla de contenidos}
\fi
\ifdefined\listfigurename
  \renewcommand*\listfigurename{Listado de Figuras}
\else
  \newcommand\listfigurename{Listado de Figuras}
\fi
\ifdefined\listtablename
  \renewcommand*\listtablename{Listado de Tablas}
\else
  \newcommand\listtablename{Listado de Tablas}
\fi
\ifdefined\figurename
  \renewcommand*\figurename{Figura}
\else
  \newcommand\figurename{Figura}
\fi
\ifdefined\tablename
  \renewcommand*\tablename{Tabla}
\else
  \newcommand\tablename{Tabla}
\fi
}
\makeatother
\usepackage{bookmark}
\IfFileExists{xurl.sty}{\usepackage{xurl}}{} % add URL line breaks if available
\urlstyle{same}
\hypersetup{
  pdftitle={Práctica 2: Un entorno de programación en Java},
  pdflang={es},
  colorlinks=true,
  linkcolor={uzaBlueDark},
  filecolor={Maroon},
  citecolor={Blue},
  urlcolor={uzaBlue},
  pdfcreator={LaTeX via pandoc}}

\title{Práctica 2: Un entorno de programación en Java}
\author{}
\date{}
\begin{document}
\maketitle

\renewcommand*\contentsname{Tabla de contenidos}
{
\hypersetup{linkcolor=}
\setcounter{tocdepth}{2}
\tableofcontents
}

\clearpage

\section{Objetivos de la práctica}\label{objetivos-de-la-pruxe1ctica}

Los objetivos de la segunda práctica de la asignatura son los siguientes:

\begin{itemize}
\tightlist
\item
 Familiarizarse con el ordenador y los elementos básicos de su sistema operativo.
\item
 Familiarizarse con el entorno de desarrollo Java Eclipse.
\end{itemize}

\section{El entorno de desarrollo Eclipse}\label{el-entorno-de-desarrollo-eclipse}

Para facilitar el desarrollo de programas, existen entornos que permiten trabajar más cómodamente, tales como NetBeans, Eclipse o IntelliJ. Nosotros recomendamos el uso de \textbf{\emph{Eclipse}}, un entorno de desarrollo de libre distribución.

Inicia \emph{Eclipse} haciendo doble clic sobre el ejecutable \emph{Eclipse} en la carpeta \texttt{software\_disponible}, o bien buscando \emph{Eclipse} en la barra de tareas. Para instalar esta aplicación en otros ordenadores puede consultarse el \hyperref[appendix-installation]{Apéndice}. Al iniciar \emph{Eclipse} aparecerá la pantalla de bienvenida que se muestra en la Figura~\ref{fig-init_eclipse}.

\begin{figure}
\PracticeCenteredImage[keepaspectratio]{../resources/practica2/images/init_eclipse.png}
\caption{Pantalla de bienvenida.}\label{fig-init_eclipse}
\end{figure}

Para poder comenzar a trabajar con el entorno de desarrollo será necesario definir un espacio de trabajo o \emph{workspace}, tal y como se muestra en la Figura~\ref{fig-workspace_eclipse}. Simplemente se trata de una carpeta donde guardaremos nuestros proyectos. Por ejemplo, se puede establecer el directorio \texttt{C:/PracticasARQ} como directorio de trabajo.

\begin{figure}
\PracticeCenteredImage[keepaspectratio]{../resources/practica2/images/workspace_eclipse.png}
\caption{Configuración del directorio de trabajo.}\label{fig-workspace_eclipse}
\end{figure}

En los siguientes subapartados aparecen solo algunas de las características de \emph{Eclipse} o recomendaciones acerca de su uso. Si tienes más dudas sobre cómo utilizar \emph{Eclipse}, deberás consultar la ayuda en línea (opción de menú \texttt{Help\ \textgreater{}\ Help\ Contents}), o bien accediendo a su \href{http://help.eclipse.org/}{página web}, donde se encuentra también toda la documentación del programa.

\subsection{Pantalla de bienvenida}\label{pantalla-de-bienvenida}

Al abrir \emph{Eclipse}, nos aparecerá una ventana de bienvenida (Figura~\ref{fig-welcome_eclipse}) con varias opciones, de las cuales cabe destacar las siguientes:

\begin{figure}
\PracticeCenteredImage[keepaspectratio]{../resources/practica2/images/welcome_eclipse.png}
\caption{Pantalla de bienvenida.}\label{fig-welcome_eclipse}
\end{figure}

\textbf{Overview}. Proporciona un vistazo general de las características principales del entorno de desarrollo, sobre todo para familiarizarse con la interfaz de usuario.

\textbf{Tutorials}. Otra forma de conocer todas las características de Eclipse es ejecutando alguno de los tutoriales de manejo del entorno de desarrollo disponibles.

\textbf{Samples}. Una manera de conocer cómo funciona la herramienta es ejecutar algunos de los ejemplos que Eclipse pone a nuestra disposición.

\textbf{What's new}. Para los que han utilizado con anterioridad la herramienta \emph{Eclipse}, esta utilidad describe las nuevas características que proporciona la versión instalada frente a versiones anteriores.

\textbf{Workbench}. Usando esta opción se accede directamente al espacio de trabajo de Eclipse. Se encuentra arriba a la derecha, junto a un texto que dice \emph{Hide}. Esta es la opción que utilizaremos para acceder al entorno de desarrollo cada vez que iniciemos Eclipse, una vez que ya nos hayamos familiarizado con la herramienta.

Dentro de las opciones que se muestran en la pantalla de bienvendia encontramos \emph{Create a Java project}. Utilizaremos esta opción cada vez que queramos crear un nuevo proyecto, por ejemplo, para cada práctica de programación en Java de este curso. Para esta práctica, podemos crear un nuevo proyecto al que llamaremos \texttt{practica2}.

Si cerramos la pestaña de bienvenida (\emph{Welcome}, arriba a la izquierda), nos conduce al entorno de trabajo que se explica en la siguiente sección.

\subsection{Entorno de trabajo}\label{entorno-de-trabajo}

El \textbf{entorno de trabajo o \emph{workbench}} nos proporciona una interfaz de usuario intuitiva, bien estructurada y sencilla de utilizar, donde están bien definidas las distintas zonas de trabajo. En la Figura~\ref{fig-workbench_eclipse} se puede observar la interfaz de usuario de \emph{Eclipse}.

Las diferentes partes en las que está dividida la interfaz son:

\begin{itemize}
\tightlist
\item
 En la \textbf{zona izquierda}, se encuentra un menú (\textbf{\emph{Package Explorer}}) que agrupa los componentes de cada uno de los proyectos abiertos.
\item
 En la \textbf{zona central} de la pantalla se encuentra un \textbf{editor}, que sirve para ver y modificar los ficheros Java con los que trabajamos.
\item
 En la \textbf{zona de la derecha}, hay un resumen (\emph{Outline}) que nos indica la \textbf{estructura del fichero activo} en el editor (clases, métodos, variables, etc.).
\item
 En la \textbf{parte inferior del editor} se encuentran algunas utilidades que nos permiten conocer el estado del proyecto, incluyendo \textbf{errores} (pestaña \emph{Problems}) y la \textbf{salida} obtenida en los programas ejecutados (pestaña \emph{Console}).
\end{itemize}

\begin{figure}
\PracticeCenteredImage[keepaspectratio]{../resources/practica2/images/workbench_eclipse.png}
\caption{Espacio de trabajo de \emph{Eclipse IDE}.}\label{fig-workbench_eclipse}
\end{figure}

\subsection{Crear un proyecto}\label{crear-un-proyecto}

Para la realización de cada práctica se recomienda \textbf{crear un proyecto nuevo}. De esta forma se pueden agrupar los ficheros correspondientes a una determinada práctica dentro de un proyecto para compilar y ejecutar programas individuales con una configuración concreta.

Los proyectos se crean a través de la opción \texttt{File\ \textgreater{}\ New\ \textgreater{}\ Java\ Project}. Para crear el nuevo proyecto hay que dotarlo de un nombre identificativo (por ejemplo, \texttt{practica2}). Dicho proyecto se guarda en un directorio del equipo, dentro del \emph{workspace} seleccionado al lanzar \emph{Eclipse}. Conviene seleccionar la opción \textbf{\emph{Create separate source and output folders}}. En la Figura~\ref{fig-new_project} se pueden observar todas las propiedades que se deben definir en la creación del proyecto.

Teniendo en cuenta que nuestro entorno de trabajo se encuentra en \texttt{C:\textbackslash{}practicasARQ}, el nuevo proyecto que has creado se encuentra en la carpeta:

\texttt{C:\textbackslash{}practicasARQ\textbackslash{}practica2}

Abre dicha carpeta en el explorador de archivos del sistema operativo y observa que se han creado dos subcarpetas. En \texttt{src} se almacenarán los ficheros con el código en lenguaje Java (\emph{source}), mientras que en la carpeta \texttt{bin} se almacenan los ficheros binarios (\emph{binaries}).

\begin{figure}
\PracticeCenteredImage[width=0.7\linewidth,height=\textheight,keepaspectratio]{../resources/practica2/images/new_project.png}
\caption{Propiedades que debemos configurar al crear un nuevo proyecto.}\label{fig-new_project}
\end{figure}

\begin{tcolorbox}[enhanced jigsaw, toprule=.15mm, coltitle=black, titlerule=0mm, opacitybacktitle=0.6, breakable, colback=white, colframe=quarto-callout-warning-color-frame, colbacktitle=quarto-callout-warning-color!10!white, opacityback=0, title={Configuración del proyecto}, bottomtitle=1mm, rightrule=.15mm, toptitle=1mm, left=2mm, bottomrule=.15mm, leftrule=.75mm, arc=.35mm]

Recuerda \textbf{desmarcar la opción \emph{Create module-info.java}} y \textbf{marcar la opción \emph{Create separate source and output folders}} para que se cree una carpeta para el código fuente y otra para los ficheros binarios.

\end{tcolorbox}

\subsection{Crear una aplicación}\label{crear-una-aplicaciuxf3n}

Los lenguajes orientados a objetos, como Java, giran en torno al concepto de clase. Todas las variables y métodos de Java deben pertenecer a una \textbf{clase}. Una clase es una colección de datos (variables) y de métodos (funciones) que operan sobre dichos datos. Cada clase pública debe ir en un fichero con extensión \texttt{.java}, cuyo nombre es exactamente igual que el de la clase. Por ejemplo, en un fichero \texttt{HolaMundo.java}, definiremos una clase \emph{HolaMundo}.

\begin{figure}
\PracticeCenteredImage[width=0.75\linewidth,height=\textheight,keepaspectratio]{../resources/practica2/images/class_skeleton.png}
\caption{Ejemplo de estructura de una clase Java.}\label{fig-class_skeleton}
\end{figure}

Para \textbf{añadir nuevas clases} al proyecto utilizaremos la opción de menú \texttt{File\ \textgreater{}\ New\ \textgreater{}\ Class}. Selecciona \texttt{src} como subcarpeta e introduce \texttt{HolaMundo} como nombre del archivo que contendrá el programa (clase) que deberás escribir a continuación. La opción \texttt{public\ static\ void\ main(String{[}{]}\ args)} debe marcarse para que la clase sea ejecutable, ejecutándose por tanto el método \emph{main} de la clase. La Figura~\ref{fig-new_class} muestra la configuración de nuestra nueva clase Java.

\begin{figure}
\PracticeCenteredImage[width=0.6\linewidth,height=\textheight,keepaspectratio]{../resources/practica2/images/new_class.png}
\caption{Creación de una nueva clase de Java en \emph{Eclipse IDE}.}\label{fig-new_class}
\end{figure}

Fijaos en que \textbf{no hemos indicado el nombre de un paquete, por lo que la clase formará parte un paquete por defecto (\emph{default})}. En futuras sesiones de prácticas hablaremos de qué es un paquete, y cómo crearlo para estructurar las clases que creemos en cada proyecto.

Al finalizar la creación, se generará automáticamente un pequeño esqueleto con el código inicial común a cada clase, que deberá completarse para contener el siguiente código:

\PracticeCode{../resources/practica2/snippets/java/hola-mundo-errors.java}

\begin{tcolorbox}[enhanced jigsaw, toprule=.15mm, coltitle=black, titlerule=0mm, opacitybacktitle=0.6, breakable, colback=white, colframe=quarto-callout-warning-color-frame, colbacktitle=quarto-callout-warning-color!10!white, opacityback=0, title={Errores de compilación}, bottomtitle=1mm, rightrule=.15mm, toptitle=1mm, left=2mm, bottomrule=.15mm, leftrule=.75mm, arc=.35mm]

El código anterior contiene algunos \textbf{errores introducidos intencionadamente} para que se produzcan \textbf{errores de compilación}.

\end{tcolorbox}

El editor de \emph{Eclipse} es un editor típico del entorno Windows: se puede copiar, cortar, pegar, buscar y sustituir. Tiene un sistema de tabulación inteligente que facilita la tarea de escritura. Se puede conseguir aumentar o disminuir la \textbf{tabulación de una porción de código seleccionando} dicha porción y pulsando tabulador (\texttt{TAB}) o mayúsculas-tabulador (\texttt{Shift\ +\ TAB}), respectivamente. Además, el editor de \emph{Eclipse} tiene un sistema de \textbf{autocompletado} que se activa pulsando \texttt{Ctrl\ +\ Space} y que facilita la escritura de código, por ejemplo, para completar el nombre de una clase o método a partir de las primeras letras escritas.

\textbf{Guarda los cambios} que hemos introducido pulsando los iconos con disquetes que hay en la barra de herramientas superior, pulsando \texttt{Ctrl-S} o bien con \texttt{File\ \textgreater{}\ Save}.

\subsection{Compilación de aplicaciones}\label{compilaciuxf3n-de-aplicaciones}

Para compilar un programa, es necesario \textbf{traducir el código fuente escrito en lenguaje Java a un código binario} que pueda ser interpretado por la máquina virtual de Java. En el entorno \emph{Eclipse}, la compilación se realiza \textbf{automáticamente cada vez que se guarda un fichero}, aunque también es posible compilar manualmente utilizando la opción de menú \texttt{Project\ \textgreater{}\ Build\ Project}. Si el programa contiene errores de compilación, \emph{Eclipse} nos lo indicará en la zona inferior del editor, en la subpestaña \textbf{\emph{Problems}} (ver Figura~\ref{fig-compilation_errors}). En esta subpestaña se nos informará sobre el proceso de compilación indicando los posibles errores. \textbf{Puedes seleccionar cada error mediante doble click, y se indicará dentro del código dónde se encuentra el error}.

\begin{figure}
\PracticeCenteredImage[keepaspectratio]{../resources/practica2/images/file_types.png}
\caption{Proceso de ejecución de un programa.}\label{fig-file_types}
\end{figure}

\begin{figure}
\PracticeCenteredImage[keepaspectratio]{../resources/practica2/images/compilation_errors.png}
\caption{Compilación de un programa con errores.}\label{fig-compilation_errors}
\end{figure}

Si trabajásemos sin un entorno de desarrollo, sería equivalente a abrir una ventana de órdenes del sistema operativo y teclear:

\PracticeCode{../resources/practica2/snippets/commands/compile-hola-mundo.txt}

\begin{tcolorbox}[enhanced jigsaw, toprule=.15mm, coltitle=black, titlerule=0mm, opacitybacktitle=0.6, breakable, colback=white, colframe=quarto-callout-note-color-frame, colbacktitle=quarto-callout-note-color!10!white, opacityback=0, title={Resolución de errores de compilación}, bottomtitle=1mm, rightrule=.15mm, toptitle=1mm, left=2mm, bottomrule=.15mm, leftrule=.75mm, arc=.35mm]

El primer error se produce porque el nombre correcto del método es \texttt{nextLine}; el segundo porque se ha escrito mal la variable \texttt{nombre}. Corrige ambos errores y recompila guardando el fichero (\texttt{Ctrl-S}) o utilizando la opción \texttt{Project\ \textgreater{}\ Build\ project}. Ahora no te deberá mostrar ningún error. El fichero compilado se llama \texttt{HolaMundo.class} y se habrá creado en la subcarpeta \texttt{bin}.

\end{tcolorbox}

\subsection{Ejecución de aplicaciones}\label{ejecuciuxf3n-de-aplicaciones}

Una vez que el programa se ha compilado correctamente, es posible ejecutarlo. Para ello, en el entorno Eclipse se utiliza la opción de menú \textbf{\texttt{Run\ \textgreater{}\ Run\ as:\ Java\ Application}} o bien el icono de \textbf{ejecutar} en la barra de herramientas superior. Si el programa se ejecuta correctamente, aparecerá una nueva pestaña llamada \emph{Console} en la parte inferior del editor, donde se mostrará la salida del programa.

\begin{tcolorbox}[enhanced jigsaw, toprule=.15mm, coltitle=black, titlerule=0mm, opacitybacktitle=0.6, breakable, colback=white, colframe=quarto-callout-note-color-frame, colbacktitle=quarto-callout-note-color!10!white, opacityback=0, title={Ejecución sin un entorno de desarrollo}, bottomtitle=1mm, rightrule=.15mm, toptitle=1mm, left=2mm, bottomrule=.15mm, leftrule=.75mm, arc=.35mm]

Si trabajásemos directamente con el JDK deberíamos invocar al intérprete (comando \texttt{java}) pasándole como parámetro el fichero compilado sin poner la extensión \texttt{.class}. Así pues, en el intérprete de órdenes teclearíamos:

\PracticeCode{../resources/practica2/snippets/commands/run-hola-mundo.txt}

\end{tcolorbox}

\subsection{Depuración de errores}\label{depuraciuxf3n-de-errores}

En el proceso de desarrollo de un programa, es habitual que se produzcan errores de ejecución. Para solucionar estos errores, es necesario \textbf{depurar el programa}, es decir, \textbf{ejecutar el programa paso a paso} para observar el estado de las variables y comprender mejor el funcionamiento real del programa (ver Figura~\ref{fig-debug_eclipse}).

\begin{tcolorbox}[enhanced jigsaw, toprule=.15mm, coltitle=black, titlerule=0mm, opacitybacktitle=0.6, breakable, colback=white, colframe=quarto-callout-note-color-frame, colbacktitle=quarto-callout-note-color!10!white, opacityback=0, title={Ejecución paso a paso}, bottomtitle=1mm, rightrule=.15mm, toptitle=1mm, left=2mm, bottomrule=.15mm, leftrule=.75mm, arc=.35mm]

La opción \texttt{Step\ into} ejecuta el código \emph{instrucción a instrucción}, mientras que \texttt{Step\ over} \emph{ejecuta las llamadas a métodos como sentencias simples} (sin entrar a ejecutar paso a paso cada instrucción dentro del método).

\end{tcolorbox}

\begin{figure}
\PracticeCenteredImage[keepaspectratio]{../resources/practica2/images/debugging.png}
\caption{Depuración de un programa en Eclipse.}\label{fig-debug_eclipse}
\end{figure}

Para ello, seguimos los siguientes pasos:

\begin{itemize}
\tightlist
\item
 Utiliza la opción de menú \texttt{Run\ \textgreater{}\ Toggle\ line\ breakpoint} para establecer (o eliminar) puntos de parada sobre el código que queremos ejecutar paso a paso. La ejecución paso a paso comenzará en la línea que contenga el punto de ruptura. También se puede utilizar el menú desplegable que aparece al pulsar el botón derecho del ratón sobre el margen izquierdo de la ventana que muestra el código de una clase Java, justo al lado del número de línea.
\item
 Utiliza la opción \texttt{Run\ \textgreater{}\ Debug}, o \texttt{Debug} en la barra de herramientas superior, para lanzar la ejecución en modo depuración de errores. Aparecerá una ventana de \emph{debugger} (depurador) con el estado de todas las variables.
\item
 Utiliza \texttt{Run\ \textgreater{}\ Step\ into\ (F5)} y \texttt{Run\ \textgreater{}\ Step\ over\ (F6)} para ejecutar el código instrucción a instrucción. \texttt{Step\ over} ejecuta las llamadas a métodos como sentencias simples (sin entrar a ejecutar paso a paso cada instrucción dentro del método).
\end{itemize}

Para practicar todo lo visto hasta ahora, haz lo siguiente:

\begin{itemize}
\tightlist
\item
 Crea una nueva clase llamada \textbf{Ejercicio} (comenzando con una letra mayúscula).
\item
 Pega el siguiente código en la clase creada.
\end{itemize}

\PracticeCode{../resources/practica2/snippets/java/ejercicio-errors.java}

\begin{itemize}
\tightlist
\item
 Soluciona todos los errores de compilación en el programa anterior.
\item
 Ejecuta el programa paso a paso para observar cuál es el valor que toma la variable \texttt{x} justo antes de finalizar el programa.
\end{itemize}

\subsection{Exportar proyectos}\label{exportar-proyectos}

Las opciones de \textbf{exportar e importar proyectos} sirven para \textbf{transferir proyectos de un ordenador a otro}. Por ejemplo, podemos exportar un proyecto \emph{Eclipse} a una carpeta de una memoria USB y, en otro ordenador, importar la carpeta.

Para exportar un proyecto de \emph{Eclipse} a una carpeta de ficheros, haremos lo siguiente:

\begin{itemize}
\tightlist
\item
 Utiliza la opción de menú \texttt{File\ \textgreater{}\ Export} y elegir \texttt{General\ \textgreater{}\ File\ system} (ver Figura~\ref{fig-export_project}).
\end{itemize}

\begin{figure}
\PracticeCenteredImage[keepaspectratio]{../resources/practica2/images/export_workspace.png}
\caption{Exportación de proyectos a una carpeta de ficheros.}\label{fig-export_project}
\end{figure}

\begin{itemize}
\tightlist
\item
 Marcaremos el proyecto a exportar y todos sus elementos (carpetas en la parte izquierda, ficheros concretos en la parte derecha). En general, sólo se suele exportar la carpeta \texttt{src}, dado que los ficheros binarios (\emph{bin}) se generan a partir de la compilación.
\item
 En la zona \texttt{To\ directory} se indicará el nombre de la carpeta que vamos a crear.
\item
 Pulsa \texttt{Finish} para finalizar.
\end{itemize}

Si ha ido todo bien, habrás generado una carpeta con la estructura de subcarpetas del proyecto actual. Existen otras maneras de exportar proyectos, pero esta es la más sencilla. Por otro lado, no es estrictamente necesario exportar el proyecto, sino que es suficiente con guardar los ficheros \texttt{.java} ubicados en la carpeta \texttt{src}.

\subsection{Importar proyectos}\label{importar-proyectos}

Por otro lado, para importar un proyecto tenemos varias alternativas:

\begin{itemize}
\tightlist
\item
 Copiamos la carpeta de ficheros exportada en el \emph{workspace} del ordenador. Si existe otra carpeta con una versión previa del proyecto (por ejemplo, si estamos practicando la exportación y la importación en el mismo ordenador), deberemos borrar la versión previa para poder importarla correctamente.
\item
 Utilizamos la opción \texttt{File\ \textgreater{}\ Import} y a continuación seleccionamos \texttt{General\ \textgreater{}\ Existing\ projects\ into\ workspace} (ver Figura~\ref{fig-import_project}). De esta manera, se importará el proyecto completo, incluyendo la configuración de compilación y ejecución. Esta es la opción recomendada para importar proyectos a \emph{Eclipse}.
\end{itemize}

\begin{figure}
\PracticeCenteredImage[keepaspectratio]{../resources/practica2/images/import_workspace.png}
\caption{Importación de proyectos en un \emph{workspace}.}\label{fig-import_project}
\end{figure}

\begin{itemize}
\tightlist
\item
 Crear un nuevo proyecto a partir de clases Java existentes utilizando la opción \texttt{Import\ \ \textgreater{}\ General\ \textgreater{}\ File\ system}. De esta manera, podríamos importar únicamente los ficheros \texttt{.java} sin necesidad de considerar el resto del proyecto.
\end{itemize}

Por ahora recomendamos utilizar la opción \textbf{\texttt{General\ \textgreater{}\ Existing\ projects\ into\ workspace}}.

\subsection{Crear biblioteca}\label{crear-biblioteca}

Cuando se quiere distribuir una aplicación desarrollada en Java que incluye un número considerable de clases y paquetes, conviene generar una \textbf{biblioteca}. Las bibliotecas en Java son archivos con extensión \textbf{\texttt{.jar}} que comprimen, de manera similar a un \texttt{.zip}, un conjunto de clases ya compiladas (ficheros \texttt{.class}) y organizadas en paquetes. Se puede explorar el contenido de un fichero \texttt{.jar} con programas de compresión de ficheros como \emph{WinZip}.

\begin{figure}
\PracticeCenteredImage[keepaspectratio]{../resources/practica2/images/import_export_scheme.png}
\caption{Ejemplo de usos de importación y exportación de bibliotecas.}\label{fig-import_export_scheme}
\end{figure}

Para crear una biblioteca en el entorno \emph{Eclipse} debemos seleccionar el proyecto que queremos exportar, utilizar la opción de menú \texttt{File\ \textgreater{}\ Export}, seleccionar \texttt{Java\ \textgreater{}\ Runnable\ JAR\ File} y seguir los siguientes pasos, ilustrados en la Ilustración 12:

\begin{itemize}
\tightlist
\item
 En \texttt{Launch\ configuration} se elige la clase principal, es decir, la que tiene el método \texttt{main} que se ejecutará una vez que ejecutemos el fichero \texttt{.jar}.
\item
 En \texttt{Export\ destination} se indica el nombre del fichero que vamos a crear y la ruta donde se va a almacenar dicho fichero.
\item
 Pulsa \texttt{Finish} para finalizar. Si ha ido todo bien, se habrá generado un fichero \texttt{.jar}.
\end{itemize}

\begin{figure}
\PracticeCenteredImage[keepaspectratio]{../resources/practica2/images/jar.png}
\caption{Creación de una biblioteca .jar a partir de un proyecto.}\label{fig-create_jar}
\end{figure}

\begin{tcolorbox}[enhanced jigsaw, toprule=.15mm, coltitle=black, titlerule=0mm, opacitybacktitle=0.6, breakable, colback=white, colframe=quarto-callout-note-color-frame, colbacktitle=quarto-callout-note-color!10!white, opacityback=0, title={Ejecución de un programa a partir de un fichero .jar}, bottomtitle=1mm, rightrule=.15mm, toptitle=1mm, left=2mm, bottomrule=.15mm, leftrule=.75mm, arc=.35mm]

Si el programa se llama \texttt{HolaMundo.jar}, podrás ejecutarlo escribiendo en una terminal del sistema operativo lo siguiente:

\PracticeCode{../resources/practica2/snippets/commands/run-jar-hola-mundo.txt}

Para abrir una terminal en \emph{Windows}, hay que utilizar la barra de búsqueda e introducir \texttt{Símbolo\ del\ sistema}. Antes de poder ejecutar el \texttt{.jar} con la instrucción anterior, debemos posicionarnos en la carpeta que lo contenga. Para simplificar, guarda el fichero \texttt{.jar} en la carpeta raíz del disco duro, \texttt{C:}, y posteriormente escribe en la ventana de órdenes la siguiente instrucción:

\PracticeCode{../resources/practica2/snippets/commands/cd-root.txt}

\end{tcolorbox}

\subsection{Añadir bibliotecas a un proyecto}\label{auxf1adir-bibliotecas-a-un-proyecto}

Si queremos utilizar alguna de las clases ya definidas en Java, simplemente es necesario añadir una sentencia en nuestro programa para \textbf{importar la biblioteca} correspondiente. Por ejemplo, para utilizar las clases definidas en el paquete \texttt{java.util}, hemos añadido en nuestro programa la siguiente instrucción:

\PracticeCode{../resources/practica2/snippets/java/import-java-util.java}

El proceso de incluir bibliotecas diferentes a las incluidas por defecto en Java puede no ser trivial, dado que hay varias maneras; entre ellas, las siguientes:

\begin{itemize}
\tightlist
\item
 Si tenemos un fichero \texttt{.jar} con las clases compiladas, podemos añadirlo directamente.
\item
 Si tenemos las clases compiladas pero no un fichero \texttt{.jar}, nos puede interesar incluir el directorio con clases compiladas.
\item
 Si vamos a utilizar una biblioteca común, es recomendable crear una variable e incluirla como biblioteca.
\end{itemize}

En este apartado vamos a describir el primer escenario. Para ello, vamos a importar a modo de ejemplo la biblioteca \texttt{.jar} incluida en Moodle, \texttt{biblioteca.jar}.

\href{../resources/practica2/binaries/biblioteca.jar}{Descargar Fichero JAR}

Antes de ello, vamos a escribir en nuestro programa la siguiente instrucción:

\PracticeCode{../resources/practica2/snippets/java/practicas-escribe-mensaje.java}

Al compilar, se nos informa de un error: \textbf{la clase \emph{Practicas} no ha sido definida todavía}. Para solucionarlo, vamos a importar una biblioteca donde se define dicha clase. Debemos hacer click derecho sobre el nombre del proyecto y seleccionar \texttt{Properties\ \textgreater{}\ Java\ Build\ Path}. Obtendremos una ventana como la de la Figura~\ref{fig-add_jar}, donde se pueden configurar las bibliotecas que se van a utilizar en el proyecto.

A continuación, hacemos click sobre \texttt{Add\ External\ JARs} y seleccionamos el fichero a importar. Una vez hecho esto, podemos comprobar que es posible ejecutar nuestro programa sin ningún tipo de problemas, pues el código necesario se obtiene en el fichero importado.

\begin{figure}
\PracticeCenteredImage[keepaspectratio]{../resources/practica2/images/add_jar.png}
\caption{Añadir una biblioteca.jar a un proyecto.}\label{fig-add_jar}
\end{figure}

\subsection{Generar documentación}\label{generar-documentaciuxf3n}

Para generar la documentación de las clases y paquetes desarrollados se utiliza la herramienta \textbf{\emph{javadoc}}, incluida en la JDK. Desde el entorno de \emph{Eclipse}, se facilita la utilización de esta herramienta a través de la opción de menú \texttt{Project\ \textgreater{}\ Generate\ Javadoc} (Figura~\ref{fig-javadoc}). Al seleccionar esta opción, se muestra una ventana donde se deben configurar algunos parámetros para generar la documentación.

En primer lugar, seleccionamos la ubicación del ejecutable \emph{javadoc} con el botón \texttt{Configure}. Normalmente, \emph{Eclipse} detecta automáticamente su ubicación y rellena este campo por nosotros, por lo que no necesitamos hacer nada.

A continuación, se debe configurar el directorio de salida de la documentación utilizando la caja de texto \texttt{Destination} y el botón \texttt{Browse} asociado. Por último, se pulsará el botón \texttt{Finish} para finalizar y generar la documentación. La documentación podrá abrirse con cualquier navegador HTML. Pruébalo para la clase \texttt{HolaMundo}.

\begin{figure}
\PracticeCenteredImage[keepaspectratio]{../resources/practica2/images/javadoc.png}
\caption{Generación de documentación con la herramienta javadoc.}\label{fig-javadoc}
\end{figure}

\subsection{Acceso a la documentación del JDK}\label{acceso-a-la-documentaciuxf3n-del-jdk}

Java proporciona una amplia y detallada documentación acerca de las \textbf{herramientas que incluye el JDK}, así como la especificación de la API (\emph{Application Programming Interface}), es decir, la documentación de las bibliotecas disponibles por defecto en Java. Esta documentación se encuentra la página oficial de \emph{Oracle} (ver Figura~\ref{fig-jdk_doc}), y lo mismo sucede para la documentación de la API (ver Figura~\ref{fig-api_doc}). Esta documentación puede descargarse para poder consultarla sin necesidad de acceder a Internet; el material descargado puede visualizarse con cualquier navegador web.

Modifica la clase \texttt{Ejercicio} para que cada vez que se ejecute escriba por pantalla un número diferente. Para ello, modifica la instrucción \texttt{x\ =\ 1} con métodos de la clase \texttt{Math} que permitan obtener un número entero aleatorio: primero se obtendrá un número real aleatorio y luego se redondeará a un número entero. Puede buscarse dicho método en la documentación de la clase \href{https://docs.oracle.com/en/java/javase/25/docs/api/java.base/java/lang/Math.html}{\texttt{java.lang.Math}}.

\PracticeCode{../resources/practica2/snippets/java/math-random-placeholder.java}

También, modifica la última instrucción del programa para que imprima el valor de \texttt{x} con un mensaje descriptivo, por ejemplo:

\PracticeCode{../resources/practica2/snippets/java/fibonacci-print.java}

\begin{figure}
\PracticeCenteredImage[keepaspectratio]{../resources/practica2/images/jdk_doc.png}
\caption{Acceso a la documentación del JDK.}\label{fig-jdk_doc}
\end{figure}

\begin{figure}
\PracticeCenteredImage[keepaspectratio]{../resources/practica2/images/api_doc.png}
\caption{Acceso a la documentación de la API de Java.}\label{fig-api_doc}
\end{figure}

\section{Apéndice. Instalación de Eclipse y JDK}\label{appendix-installation}

En este apéndice explicaremos el proceso de instalación del software necesario para realizar las prácticas de Informática en los ordenadores personales. Existen al menos dos maneras de instalar el software:

\subsection{Instalación de Eclipse mediante un instalador}\label{instalaciuxf3n-de-eclipse-mediante-un-instalador}

\begin{itemize}
\tightlist
\item
 Visita la página web oficial de \href{https://www.eclipse.org/}{\emph{Eclipse}}.
\item
 Localiza el botón de \textbf{\emph{Download}}. Nos llevará a la página de descargas.
\item
 Localiza el panel de \emph{Eclipse Installer} (Figura~\ref{fig-installer}). Debemos reconocer nuestro sistema operativo (\emph{Windows}, \emph{Linux} o \emph{macOS}), y pulsar sobre \texttt{x86\_64}.
\item
 Una vez descargado el instalador, lo ejecutamos.
\end{itemize}

El instalador permite obtener múltiples herramientas de Eclipse, aunque nosotros sólo necesitamos el entorno de desarrollo para Java. La instalación es bastante automática, y simplemente tendremos que seleccionar Eclipse IDE for Java Developers (Figura~\ref{fig-suite_tools}) y pulsar el botón \texttt{Install} para iniciar la instalación.

\begin{figure}
\PracticeCenteredImage[keepaspectratio]{../resources/practica2/images/eclipse_suite_installation.png}
\caption{Panel de descarga del instalador de Eclipse.}\label{fig-installer}
\end{figure}

\begin{figure}
\PracticeCenteredImage[width=0.6\linewidth,height=\textheight,keepaspectratio]{../resources/practica2/images/eclipse_suite.png}
\caption{Herramientas disponibles en la suite de Eclipse.}\label{fig-suite_tools}
\end{figure}

Una vez seleccionada dicha opción, se nos dirige a una pantalla como la que se muestra en la Figura~\ref{fig-jre_installation}. El objetivo es configurar dos rutas: la localización de las herramientas de desarrollo para Java y la ruta de instalación de Eclipse. Podéis personalizar ambas rutas, pero no es recomendable hacerlo. En particular, la primera ruta apunta a un \emph{Java Development Kit} (JDK), que, si no queremos descargar junto con Eclipse, deberá instalarse por separado. Esta última opción, aunque es menos sencilla, se describe en el siguiente punto y, entre otras ventajas, permite probar los comandos java y javac en la terminal del ordenador.

\begin{figure}
\PracticeCenteredImage[width=0.75\linewidth,height=\textheight,keepaspectratio]{../resources/practica2/images/jre_installation.png}
\caption{Instalación de Eclipse para Java, indicando que queremos descargar un JRE.}\label{fig-jre_installation}
\end{figure}

Una opción alternativa consiste en indicar la ruta de un \emph{Java Runtime Environment} (JRE) (versión 21 o superior) previamente instalado (ver Figura~\ref{fig-jdk_installation}). Para la descarga e instalación del JDK de Java, por favor, dirígete a la sección B.1 del Apéndice. Por ejemplo, en la siguiente imagen, se ha detectado una instalación en la ruta \texttt{C:/Program\ Files/Java/jdk-25}, que es la ruta de instalación por defecto de un JDK.

\begin{figure}
\PracticeCenteredImage[width=0.75\linewidth,height=\textheight,keepaspectratio]{../resources/practica2/images/jdk_installation.png}
\caption{Instalación de Eclipse para Java, indicando que queremos utilizar un JDK previamente instalado.}\label{fig-jdk_installation}
\end{figure}

\subsection{Instalación alternativa de Eclipse sin instalador}\label{instalaciuxf3n-alternativa-de-eclipse-sin-instalador}

\subsubsection{Instalación del Java Software Development Kit (JDK)}\label{instalaciuxf3n-del-java-software-development-kit-jdk}

Para comenzar, debemos descargar la última versión del JDK de la página web de Oracle:

\begin{itemize}
\tightlist
\item
 Descargamos el JDK de Java de la página oficial de \href{https://www.oracle.com/es/java/technologies/downloads/}{Oracle}. En este momento, la última versión disponible es la \textbf{25}.
\item
 Debemos seleccionar el sistema operativo de nuestro ordenador personal (\emph{Linux}, \emph{macOS} o \emph{Windows}) y descargar un instalador. En \emph{Windows}, podemos descargar tanto \texttt{x64\ Installer} como \texttt{x64\ MSI\ Installer}.
\item
 Una vez descargado, abrimos el archivo para comenzar el proceso de instalación. El proceso es completamente automático y no habrá que hacer nada, salvo utilizar el botón \texttt{Next} para avanzar al siguiente paso. La ruta de instalación por defecto en Windows es \texttt{C:/Program\ Files/Java/jdk-25}. \textbf{Se recomienda no modificar esta ruta, salvo que no tengamos espacio en \texttt{C:}}.
\end{itemize}

\begin{figure}
\PracticeCenteredImage[keepaspectratio]{../resources/practica2/images/jdk_installation_standalone.png}
\caption{Paso intermedio de la instalación del JDK. Se muestra la ruta de instalación por defecto, la cual no modificaremos.}\label{fig-jdk_installation2}
\end{figure}

\subsubsection{Instalación de Eclipse}\label{instalaciuxf3n-de-eclipse}

\begin{itemize}
\tightlist
\item
 Visita la página web oficial de \emph{Eclipse}.
\item
 Localiza el botón de \texttt{Download}; nos llevará a la página de descargas.
\end{itemize}

\begin{figure}
\PracticeCenteredImage[keepaspectratio]{../resources/practica2/images/eclipse_ide_download.png}
\caption{Panel de descarga de Eclipse sin instalador.}\label{fig-eclipse_download}
\end{figure}

\begin{itemize}
\tightlist
\item
 Partiendo de un panel similar al de la Figura~\ref{fig-eclipse_download}, reconocemos nuestro sistema operativo y pulsamos sobre \texttt{x86\_64}. Se descargará un \texttt{.zip} con una carpeta dentro llamada \texttt{eclipse} (ver Figura~\ref{fig-eclipse_files}). Podemos situar dicha carpeta en cualquier directorio de nuestra elección, por ejemplo, en la unidad \texttt{C:}. Ten en cuenta que esta opción no requiere instalación, y por tanto, accederemos siempre a dicha carpeta cuando queramos iniciar \emph{Eclipse} haciendo click sobre \texttt{eclipse.exe}. En la Figura~\ref{fig-eclipse_files} se muestra el contenido del fichero \texttt{.zip} descomprimido.
\end{itemize}

\begin{figure}
\PracticeCenteredImage[keepaspectratio]{../resources/practica2/images/filesystem.png}
\caption{Ficheros contenidos en el directorio eclipse, tras descomprimir el fichero .zip descargado de la página oficial de Eclipse.}\label{fig-eclipse_files}
\end{figure}

\subsection{Pantalla de bienvenida}\label{pantalla-de-bienvenida-1}

La primera vez que se utiliza \emph{Eclipse}, aparecerá un mensaje preguntándonos si queremos excluir \emph{Eclipse} del análisis del antivirus de Windows. Como aparece en la Figura~\ref{fig-antivirus_eclipse}, podemos indicar que se excluya \emph{Eclipse} de dichos análisis para evitar ralentizar la aplicación.

\begin{figure}
\PracticeCenteredImage[keepaspectratio]{../resources/practica2/images/antivirus.png}
\caption{Pregunta de exclusión de Eclipse IDE respecto del sistema antivirus de Windows.}\label{fig-antivirus_eclipse}
\end{figure}

\section{Entrega de la solución}\label{entrega-de-la-soluciuxf3n}

\textbf{Sólo debe realizar la entrega un miembro del grupo}. Sube a Moodle el fichero \texttt{Ejercicio.java} con el código de la clase \texttt{Ejercicio}. No es necesario subir el proyecto completo ni ningún otro fichero.

\end{document}
